% =============================================================================
% cap09.tex — CONSIDERACOES FINAIS
% =============================================================================

\chapter{CONSIDERAÇÕES FINAIS}\label{cap:consideracoes}

\section{QUANTO AOS OBJETIVOS}\label{sec:quanto-objetivos}

O objetivo geral enunciado na \autoref{sec:objetivo-geral} foi alcançado. A
estrutura do galpão em madeira lamelada colada foi dimensionada e verificada
segundo a EN 1995-1-1 \cite{EN1995-1-1:2004}, por meio de programa de cálculo
automatizado, e o memorial resultante consta dos capítulos
\autoref{cap:dimensionamento} e \autoref{cap:detalhamento}. O conjunto de
quarenta verificações é atendido, com taxa de utilização máxima de 0,7902.

Os objetivos específicos foram cumpridos nos termos seguintes. A caracterização
do material e das exigências de durabilidade consta da
\autoref{sec:materiais}. As ações e as combinações constam da
\autoref{sec:acoes} e da \autoref{sec:combinacoes}. A geometria foi estabelecida
por varredura paramétrica de 1.391 candidatas, apresentada na
\autoref{sec:varredura-global}. As verificações de estado-limite último dos
elementos, dos sistemas de contraventamento e da estabilidade global constam do
\autoref{cap:dimensionamento}, e as de serviço, da \autoref{sec:verif-els}. As
ligações e a ancoragem constam do \autoref{cap:detalhamento}. O programa,
submetido a suíte de 558 testes automatizados, é reproduzido no
\autoref{ap:codigo-fonte}.

\section{PRINCIPAIS RESULTADOS}\label{sec:principais-resultados}

Sete resultados merecem destaque, por contrariarem expectativas correntes ou por
terem determinado decisões de projeto.

O primeiro é o modelo de apoio horizontal. A verificação da rigidez relativa
entre a viga, que se opõe ao afastamento com 208.462 kN/m, e o pilar que a
apoia, que oferece rigidez nula por ser articulado, demonstrou que o empuxo não
se desenvolve. A viga trabalha como biapoiada, e o momento de projeto é 71,4\%
superior ao que o modelo de arco indicaria. A consequência favorável é que não
há esforço horizontal por efeito de arco a transmitir à fundação.

O segundo é que nenhuma verificação da viga governa o dimensionamento. A mais
solicitada, o cisalhamento, atinge 0,3814. O que governa é o contraventamento da
cobertura, com 0,7902, cuja solicitação responde à área exposta ao vento e é
praticamente insensível à seção da viga. Em uma edificação aberta e de cobertura
leve, a estabilização determina o projeto mais do que a resistência dos
elementos.

O terceiro é que a combinação governante em momento é a gravitacional, e não a
de vento. Ainda que a sucção produza carga de módulo comparável, ela atua em
sentido oposto ao peso próprio e é parcialmente cancelada. A sucção permanece
determinante para o arrancamento e para a ancoragem.

O quarto é o peso próprio. Calculado por parcelas a partir da geometria,
resultou em 0,8229 kN/m\textsuperscript{2}, cerca de metade do valor que uma
estimativa global usual atribuiria. Adotar valor único de tabela teria
sobrecarregado o modelo em aproximadamente oitenta por cento e mascarado a
leveza própria da solução em madeira.

O quinto é a força longitudinal de vento. A parcela de atrito sobre a cobertura,
de 776,8 kN, equivale a 2,35 vezes o arrasto de bloqueio da silhueta e governa
todo o dimensionamento do contraventamento longitudinal, dos pilares de canto e
das respectivas ancoragens. A prescrição é envolvente por construção, e a ordem
de grandeza é característica dela.

O sexto é o fendilhamento nas ligações. Por ser a resultante quase perpendicular
às fibras, a taxa sem reforço atinge 0,961. O valor atende ao limite normativo,
porém excede o limite de projeto de 0,75 adotado para modos frágeis, o que
determinou a adoção de reforço perpendicular às fibras. Foi o único ponto em que
o critério de margem adicional alterou efetivamente o detalhe construtivo.

O sétimo é a parcela de cisalhamento na flecha, de 21,2\% do total. Em peças de
madeira, cuja razão entre os módulos de elasticidade e de cisalhamento é da
ordem de dezessete, essa parcela não é desprezável, e desprezá-la subestimaria a
deformação em mais de um quinto.

\section{CONTRIBUIÇÕES}\label{sec:contribuicoes}

A contribuição principal do trabalho não está em nenhum dos resultados
numéricos, e sim no procedimento que os produziu.

A automação da cadeia de verificação permitiu submeter 1.391 geometrias ao mesmo
conjunto de quarenta verificações que produz o memorial final. Essa equivalência
foi verificada: o mecanismo de varredura reproduz, sobre a geometria adotada, as
taxas do memorial sem divergência na quarta casa decimal. Não se trata de
avaliar candidatas por modelo simplificado e depois verificar a escolhida pelo
modelo completo, prática que introduz descontinuidade entre a fase de concepção
e a de verificação.

A consequência prática é que o projeto sabe onde está no espaço de soluções. A
\autoref{tab:fronteira-volume} mostra que existem geometrias aprovadas com
metade do volume de madeira da adotada, e a
\autoref{ssec:o-que-governa} explica por que a seção mais robusta foi mantida.
Sem a varredura, a margem existiria do mesmo modo, porém sem que se soubesse da
sua existência nem da sua dimensão.

A suíte de 558 testes automatizados constitui contribuição de mesma natureza.
Ela impede que uma alteração em qualquer módulo altere silenciosamente
resultados de outro, risco real em uma cadeia de verificações mutuamente
dependentes.

\section{LIMITAÇÕES}\label{sec:limitacoes}

O trabalho apresenta limitações que condicionam a aplicação direta dos seus
resultados e que se declaram a seguir.

\subsection{Dependência de anexo nacional}\label{ssec:lim-ndp}

Dezenove parâmetros de determinação nacional adotam o valor recomendado pelo
respectivo Eurocódigo, por inexistir anexo nacional brasileiro, conforme
registrado no \autoref{qua:ndp}. Entre eles figuram os coeficientes parciais dos
materiais e os limites de deformação, que incidem diretamente sobre as taxas de
utilização. Submetida sob um anexo nacional concreto, a estrutura deve ser
reverificada.

\subsection{Verificações de ancoragem incompletas}\label{ssec:lim-ancoragem}

Três modos de ruptura do concreto previstos na EN 1992-4 \cite{EN1992-4:2018}
não foram avaliados: o fendilhamento do concreto, o arrancamento lateral e a
ruptura de borda sob cortante. Os dois primeiros dependem das dimensões
definitivas do bloco e da sua armadura, que não integram o escopo. Acresce que
as resistências apresentadas decorrem de formulação normativa, e não de ensaio
do produto efetivamente especificado.

\subsection{Dados dependentes de confirmação}\label{ssec:lim-dados}

Quatro valores adotados como dados de entrada carecem de confirmação documental.
A velocidade básica do vento e o fator estatístico não estão amarrados à
isopleta e ao grupo de classificação do endereço definitivo da obra. As cargas
de impacto de veículo correspondem a prática europeia, ela própria dependente de
anexo nacional. Os pesos da telha e das instalações são valores declarados, não
extraídos de ficha técnica. A resistência do fixador da telha é valor típico de
catálogo, e governa a verificação mais solicitada da cobertura.

\subsection{Escopo não coberto}\label{ssec:lim-escopo}

A fundação não é dimensionada, por inexistir sondagem; entregam-se os esforços
por base, conforme a \autoref{tab:cargas-fundacao}. A resistência ao fogo não é
verificada. Não integram o escopo o projeto de instalações, o detalhamento de
vedações, o plano de montagem e a análise de custos.

\section{RECOMENDAÇÕES PARA TRABALHOS FUTUROS}\label{sec:recomendacoes}

Decorrem das limitações expostas, e de resultados que o trabalho produziu sem
esgotar, as seguintes linhas de continuidade.

A primeira é completar as verificações de ancoragem da
EN 1992-4 \cite{EN1992-4:2018}, em conjunto com o dimensionamento do bloco de
fundação, uma vez disponível a investigação geotécnica. É a lacuna mais
imediata, por afetar um elemento já verificado como o mais solicitado do caminho
de carga.

A segunda é a otimização conjunta da viga e dos sistemas de estabilização. A
varredura conduzida tratou a seção da viga como variável e os sistemas de
contraventamento como fixos. Como se demonstrou que estes últimos é que governam
o projeto, o ótimo do conjunto não coincide necessariamente com o obtido, e uma
varredura que inclua o número de painéis e os diâmetros das diagonais como
variáveis produziria resposta mais bem fundamentada.

A terceira é a comparação sistemática com a ABNT NBR 7190-1
\cite{NBR7190-1:2022}. A norma brasileira, revista em 2022, aproximou-se do
formato dos Eurocódigos, e uma verificação paralela permitiria quantificar a
divergência entre os dois conjuntos normativos para esta tipologia, contribuindo
para a discussão sobre a suficiência do texto brasileiro quanto à madeira
lamelada colada.

A quarta é a validação experimental ou numérica independente do modelo de apoio
horizontal, cuja consequência sobre o momento de projeto é de 71,4\%. Ainda que
a análise de rigidez relativa seja conclusiva, a verificação por modelo de
elementos finitos ou por medição em protótipo consolidaria o resultado.

A quinta é a análise de custo. A \autoref{tab:fronteira-volume} quantifica o
volume de madeira de cada solução, porém o custo total depende também das
ligações, do aço de contraventamento e da montagem, cujos comportamentos com a
geometria são distintos. A varredura do número de painéis, na
\autoref{tab:varredura-paineis}, já sugere que a intuição de custo pode
enganar: o peso de aço mostrou-se praticamente independente do número de
painéis.
